
\chapter{SYSTEM ANALYSIS}

\ChapterSection{System Analysis}


Krishi Vaidya crop health detection project has four components which are mobile app, backend API, YOLO detector, and vision model. The initial step of this project flow is to take a picture of the crop, which would be done by farmers or field workers using the mobile app features. The mobile app supports local storage when it is not connected to the internet and synchronises the data when the internet is connected. The mobile send the captured image to the backend, which then sends the request to the VLM server, which processes the image and responds with the disease and others detail. After receiving the request from the VLM, the backend generates the advisory feedback and sends the response to the mobile app, which would present those results to the farmers or field workers. 

\ChapterSubsection{Requirement Analysis}

\subsubsection{Functional Requirements}

The different tools and system that we have used to build the Krishi Vaidya project would determine the functional requirements. We had taken these details from the implemented mobile features, the backend route groups, and the active model pipelines alongside their specific validation rules.

\begin{xltabular}{\textwidth}{@{}p{2.1cm}X@{}}
\caption{Functional requirements of Krishi Vaidya}\label{tab:functional-requirements}\\
\toprule
\textbf{ID} & \textbf{Requirement} \\
\midrule
\endfirsthead
\multicolumn{2}{l}{\footnotesize\textit{Table \thetable{} \textemdash{} continued from previous page}}\\[3pt]
\toprule
\textbf{ID} & \textbf{Requirement} \\
\midrule
\endhead
\endfoot
\bottomrule
\endlastfoot
FR-01 & The mobile application shall enable users to complete registration, language selection and configure app usage settings. \\
FR-02 & The backend shall allow users to register their crop plots, create crop records and notes, add tasks and view the scan history. \\
FR-03 & The client shall allow users to take/choose the photo of crops, see the preview photos and be allowed to make adjustments before submitting photos to the system to get the diagnoses. \\ 
FR-04 & The application backend should verify authenticated requests and verify authorized on the scan record and related resources ownership requests. \\

FR-05 & Uploaded crop images need to validate on the backend. And it immediately rejects the invalid payloads not to be included in the maiun diagnosis flow. \\
FR-06 & The VLM inference service processes the image-diagnosis requests from verified-auth users. Returns the disease prediction in our case of CLEAR-parsebile-JSON fields. \\
FR-07 & The backend then builds the structured response for the mobile app. This payload contain core diagnosis, disease severity metrics (like confidence scores)and exact treatment and prevention advisories. \\
FR-08 & Backend  stores the scan records and crop data in the database. Mobile handles tasks, ongoing notes, and the current diagnosis or advisory state. \\
FR-09 & In the rural areas of Nepal, where there is still no strong internet connection, the Krishi Vaidya mobile app is designed to work even when it is not connected to the internet. When the app is connected to the internet, the locally saved data from offline use is synchronised with the server or backend. 
 \\
\end{xltabular}
\subsubsection{Non-Functional Requirements}

\begin{xltabular}{\textwidth}{@{}p{3.0cm}X@{}}
\caption{Non-functional requirements of Krishi Vaidya}\label{tab:non-functional-requirements}\\
\toprule
\textbf{Category} & \textbf{Requirement} \\
\midrule
\endfirsthead
\multicolumn{2}{l}{\footnotesize\textit{Table \thetable{} \textemdash{} continued from previous page}}\\[3pt]
\toprule
\textbf{Category} & \textbf{Requirement} \\
\midrule
\endhead
\endfoot
\bottomrule
\endlastfoot
Usability & We had started to design the mobile app from the user's perspective to ensure that they would feel comfortable while using it and could take use of all the features that it could offer.   \\
Reliability & When designing the Krishi Vaidya project, we were aware of the cost of misdiagnosis; therefore, throughout the flow, we have implemented different levels of filters and validations that would gracefully handle errors or incorrect responses. \\
Offline support & It is common not to have internet connections in the crop fields. And this should not prevent users or farmers from using the mobile app. Therefore, we have designed the mobile app in such a way that it provides a good experience even when it's not connected to the internet. It could automatically detect the network connection and resume the tasks that were saved or pending during offline use.  \\
Security & We have used JWT for the authentication and strict request validation. And also added defenses include rate limits, locked-down CORS settings, and Helmet middleware, to secure protected route boundaries. \\
Maintainability & We have used the dependency injection design pattern so developers can test or swap out validation logic and model clients without breaking the rest of the application. \\
Performance & The YOLO model, which is used, is lightweight, and quantised variants for field analysis. VLM inference remains a heavy, server-bound operation. \\
Testability & Test suites focused on the model utilities and VLM API behaviour, wrapping backend scan services alongside mobile inference components. \\

\end{xltabular}
\ChapterSubsection{Feasibility Analysis}

\begin{xltabular}{\textwidth}{@{}p{2.7cm}X@{}}
\caption{Feasibility analysis of Krishi Vaidya}\label{tab:feasibility-analysis}\\
\toprule
\textbf{Aspect} & \textbf{Analysis} \\
\midrule
\endfirsthead
\multicolumn{2}{l}{\footnotesize\textit{Table \thetable{} \textemdash{} continued from previous page}}\\[3pt]
\toprule
\textbf{Aspect} & \textbf{Analysis} \\
\midrule
\endhead
\endfoot
\bottomrule
\endlastfoot

Technical & The centrepiece is ready; devs have built YOLO and VLM Pipelines, they are wrapped inside an Express / Mongoose Backend and run live via a React Native/ Expo mobile application. The main challenge remains reliability of diagnostics; unparseable outputs from the VLM could cause instability. \\
Operational & The phone interface has the ability to process live scans in the field, record a crop's full details and save details locally to ensure off-line working on local data. Workflow has been tested and tested on a hardware successfully. The real -world working still needs to be verified by agricultural specialists. \\
Economic & It's built out of tech 100\% accessible to student developers and free and open source. The majority of costs are funding the raw computer resources to run the model and to continue training. \\
Schedule & Development refers to CSIT project targets and deadlines. It demands presentation of an initial proposal and progress after midterm, along with final demonstration before the board examination time. This manuscript forms the core implementation evidence of the final project submission.
\end{xltabular}

\ChapterSubsection{Object-Oriented Analysis}

The architecture for the seven objects; User, CropPlot and Scan, \& Diagnosis, Advisory, CropTask \& CropNote revolves around three base objects and related ones. In backend the data is directly exposed with Mongoose models \& repository objects by the means of these seven domain objects. Mobile clients work a little differently by having mobile repositories with the backed datasource based on SQLite. State remains to be local. With the Typescript objects mobile clients can have those local states for the feature viewmodels exposed on the mobile applications.

\begin{figure}[htbp]
\centering
\includegraphics[width=\textwidth, height=0.6\textheight, keepaspectratio]{data/flow_digram_data/build/system_analysis/01_use_case_diagram.pdf}
\caption{ Diagnosis-support workflow}
\label{fig:system-use-case-diagram}
\vspace{1.5cm}
\end{figure}

\begin{figure}[htbp]
\centering
\includegraphics[width=\textwidth, height=0.6\textheight, keepaspectratio]{data/flow_digram_data/build/system_analysis/02_object_model_diagram.pdf}
\caption{Domain entities}
\label{fig:system-object-model}
\end{figure}

\begin{figure}[htbp]
\centering
\includegraphics[width=\textwidth, height=0.6\textheight, keepaspectratio]{data/flow_digram_data/build/system_analysis/03_sequence_diagram.pdf}
\caption{Sequence diagram of the online scan diagnosis workflow}
\label{fig:scan-sequence-diagram}
\vspace{0.5cm}
\end{figure}

\begin{figure}[htbp]
\centering
\includegraphics[width=\textwidth, height=0.6\textheight, keepaspectratio]{data/flow_digram_data/build/system_analysis/04_activity_diagram.pdf}
\caption{ Scan submission under online and offline conditions}
\label{fig:scan-activity-diagram}
\end{figure}